




\section{Акустовідпал $\gamma$--індукованих дефектів у структурах Au--SiO$_2$--Si}

Відомо, що радіаційні дефекти можуть ефективно взаємодіяти з пружними коливаннями.
Одним із проявів цього явища є відпал дефектів внаслідок УЗО при температурах, значно нижчих ніж це відбувається при без--звуковому нагріванні.
Подібні процеси спостерігалися у монокристалах Si \cite{OstrovRadSi,Podolian2012r,PodolHivr,YOlikh2006TPLr}, Ge \cite{Olikh:FTP1996},
напівпровідникових \cite{OlikhProc,OstrovFTTRad} та лужногалоїдних \cite{UST:OstrovCsI} сполук.
Зазвичай вони пов'язуються з розпадом радіаційно--утворених комплексів та АІ дифузією дефектів до різноманітних стоків.
Крім того, в літературі показана можливість часткового відновлення за допомогою УЗО параметрів опромінених поверхнево--бар'єрних структур, таких як, наприклад,
сонячні елементи \cite{YOlikh2007TPLr} чи світловипромінюючі діоди \cite{US:LED,UST:LED_SM}.
Найчастіше дослідження проводяться для $\gamma$--опромінених структур, проте показана можливість ефективного впливу пружних коливань і на порушення періодичності, викликані
рентгенівськими променями \cite{UST:OstrovCsI}, нейтронами \cite{Olikh:FTP1996} чи електронами \cite{US:LED,UST:LED_SM}.

   Дослідники не залишили поза увагою і можливість впливу УЗО на властивості таких промислово--важливих структур, як система
кремній --- оксид кремнію.
Зокрема, повідомлено про АІ зміни дефектного стану межі  Si--SiO$_2$ \cite{Ostap:SiO2,UST:Medvid,Zaver:2008r} та часу життя неосновних носіїв у області кремнію,
що прилягає до контакту \cite{Parchinskii2003r,Zdeb1989}.
Крім того, декілька робіт присвячено виявленню наслідків УЗО структур метал---окис---напівпровідник (МОН) на основі кремнію, опромінених $\gamma$--квантами $^{60}$Co з дозою $10^6$~рад \cite{Parchinskii2000r,Parchinskii2006r}.
Об'єктом досліджень були структури Si--SiO$_2$, виготовлені методом термічного окислення $n$--кремнію з питомим опором 0,2~Ом$\cdot$см.
%Поглинута доза опромінення складала $10^6$~рад.
У результаті вимірів ВФХ зроблено висновок про зменшення ефективного поверхневого заряду та генераційного часу життя у приконтактній області кремнію і незначне
зростання швидкості поверхневої рекомбінації.
Автори пов'язали виявлені ефекти з дифузією дефектів та розпадом домішкових асоціатів у акустичному полі, причому зазначають, що виявлені процеси
послаблені, порівняно з неопроміненими структурами, внаслідок радіаційно--стимульованої релаксації напруг на межі Si--SiO$_2$.

У підрозділі розглянуто результати досліджень впливу УЗО на перенесення заряду в опромінених структурах Au--SiO$_2$--Si.
Мета  полягала у з'ясуванні можливості відновлення працездатності діодів Шотткі, створених на основі МОН структур і суттєво деградованих внаслідок опромінення.
Хоча зазначені роботи \cite{Parchinskii2000r,Parchinskii2006r} і були певними експериментальними передумовами для наших досліджень, проте
суттєвими відмінностями представлених результатів є те, що вони отримані
а)~для систем зі значно вищою концентрацією радіаційних дефектів;
б)~при розгляді робочого режиму ДШ, тобто при проходженні струму.



Зразки для досліджень  виготовлені з кристалічного кремнію, вирощеного методом зонної плавки.
Для легування використовувалися атоми фосфору, питомий опір кристалів складав 4000~Ом$\cdot$см.
З об'ємного матеріалу було вирізано зразки у формі паралелепіпеда розмірами $1\times5\times10$~мм$^3$.
Для формування структури МОН одна з поверхонь площею 50~мм$^2$ хімічно очищувалася в розчині HF--HNO$_3$--CH$_3$COOH (об'ємне співвідношення компонент 3:5:3),
після чого в атмосферному повітрі при $T=300$~К протягом 24~год на ній формувався шар SiO$_2$.
Надалі шляхом вакуумного напилення на поверхню наносився шар золота (30~мкг/cm$^2$).
На протилежній грані за допомогою евтектики GaZn створювався омічний контакт.
Схематична структура зразків показано на рис.~\ref{figMSSi3}.

\begin{figure}[b]
\center
\includegraphics[width=0.45\textwidth]{MSSi3}%
\caption{\label{figMSSi3}
Структура зразків Au--SiO$_2$--Si
}
\end{figure}

Опромінення здійснювалося при кімнатній температурі $\gamma$--квантами $^{60}$Co, доза дорівнювала $5\cdot10^7$~рад.
Як показали вимірювання, після еквівалентного опромінення провідність монокристалічних зразків, залишаючись електронною, зменшилася приблизно в 2 рази,
що зумовлено частковою компенсацією в процесі радіаційного дефектоутворення.
Зауважимо, що питомий опір досліджених кристалів на чотири порядки більший, ніж в роботах \cite{Parchinskii2000r,Parchinskii2006r}, а відтак частка енергетичних втрат, не зумовлених іонізацією, для $\gamma$--квантів у нашому випадку значно вища.
Доза також на півтора порядку більша, і тому очікувана концентрація радіаційних дефектів суттєво вища.


УЗО опромінених структур інтенсивністю 2~Вт/см$^2$ здійснювалася за допомогою LiNbO$_3$ п'єзоелектричного перетворювача.
У зразку збуджувалися повздовжні хвилі частотою 4~МГц.
Проведено дві послідовні УЗО, кожна тривалістю 30~хв.
Під час УЗО температура зразка не перевищувала 350~К.

\begin{figure}
\center
\includegraphics[width=0.75\textwidth]{figIV_MIS}%
\caption{\label{figIV_MIS}
Прямі (а, в) та зворотні (а, б) ВАХ структур Au--SiO$_2$--Si до (криві 1)
та після (2--4) опромінення $\gamma$--квантами в напівлогарифмічному (а)
та подвійному логарифмічному (б, в) масштабах.
$t_\mathtt{UST}$ , хв: 0 (2), 30 (3), 60 (4).
$T=300$~K.
Точки --- експеримент,
лінії --- апроксимація за формулами~(\ref{eqSDIV}) (суцільні) та (\ref{eqIVTAT}) (пунктир).
Параметри апроксимацій вказані в табл.~\ref{tabMIS}.
На вставці:
залежність прямої компоненти струму неопроміненого зразка (при $V>1,6$~В) в координатах Фаулера--Нордгейма;
пряма --- лінійна апроксимація
}%
\end{figure}


На рис.~\ref{figIV_MIS} наведено ВАХ структур  Au--SiO$_2$--Si до $\gamma$--опромінення, після нього та після наступних УЗО.
З рисунка видно, що у вихідному стані ВАХ є типовою для ДШ:
при прямому зміщенні струм зумовлений ТЕ через бар'єр,
при зворотному --- визначається зміною висоти бар'єру під дією електричного поля ($I_R\sim V_R^{1/2}$).
Для апроксимації прямої гілки ВАХ був використаний вираз (\ref{eqSDIV});
результат апроксимації показано на рис.~\ref{figIV_MIS},а та б, визначені параметри
наведені в табл.~~\ref{tabMIS}.
Зауважимо, що наявність шару оксиду не дозволяє визначити ВБШ безпосередньо використовуючи величину струму насичення та формулу~(\ref{eqSDIs}), оскільки необхідно також враховувати ймовірність проходження через діелектричний прошарок \cite{OZBEK2011,Kobayashi}.
Крім того, як видно з рисунку, при $V>1$~В у неопромінених структурах величина струму перевищує значення, очікуване відповідно до виразу (\ref{eqSDIV}).
Найімовірніше причиною є тунелюванням через шар SiO$_2$, яке, загалом, може бути описане виразом (\ref{eqFowlNord}).
На користь цього припущення свідчить лінійність польової залежності величини прямого струму в координатах Фаулера--Нордгейма при великих прямих зміщеннях --- див. вставку на рис.~\ref{figIV_MIS}.
При побудові цієї залежності враховано, що величина електричного поля в прошарку оксиду товщиною $d_{ox}$ пропорційна прикладеній напрузі $F_m\sim V/d_{ox}$.






Внаслідок $\gamma$--опромінення характер проходження струму суттєво змінився --- див. криві 2 на рис.~\ref{figIV_MIS}.
Особливо це стосується зворотної гілки ВАХ (рис.~\ref{figIV_MIS},б), де кардинальні видозміни польової залежності струму свідчать про зміну механізму перенесення заряду.
При прямому зміщенні залежність $I(V)$, очікувана в рамках ТЕ моделі, спостерігається лише при напрузі, більшій 1~В (рис.~\ref{figIV_MIS},в), а величина струму суттєво менша, ніж до опромінення.
Як показала апроксимація відповідної ділянки ВАХ згідно з формулою~(\ref{eqSDIV}), у результаті опромінення відбулося значне  зростання фактора неідеальності та послідовного опору (причому зміни останнього суттєво перевищують зміни питомого опору, які спостерігаються в об'ємних зразках).
Радіаційно--індуковане збільшення $n_\mathrm{id}$, зумовлене утворенням дефектів, і є головною причиною зменшення величини термоемісійного струму.
У роботі \cite{FZSi:Rad} проведено дослідження процесів дефектоутворення у кремнії, метод вирощування і питомий опір якого збігаються із дослідженими структурами, при опромінення $\gamma$--квантами $^{60}$Co з дозою близько $9\cdot10^7$~рад.
Авторами  показано, що основними радіаційними дефектами в цьому випадку є комплекси VO$_i$, С$_i$C$_s$, $H$--центр (V$_2$O$_i$), $\Gamma$--центр та міжвузловий дефект $I^{0/-}$.
Останній є вторинним дефектом і саме з ним пов'язуються процеси компенсації (інверсії) провідності.


\begin{table}
\caption{\label{tabMIS}Визначені параметри структур Au--SiO$_2$--Si
}
\center
\begin{tabular}{|l|c|c|c|c|}
\hline
\multicolumn{5}{|c|}{Стан структури}\\\hline
$\gamma$--опромінення&$-$&$+$&$+$&$+$\\ \hline
$t_\mathtt{UST}$ , хв&0&0&30&60\\ \hline
\multicolumn{5}{|c|}{Параметр}\\\hline
$I_s$, $10^{-9}$A & 3,3& 1,1& 4,9& \\ \hline
$R_s$, $10^{4}$Ом & 1,1& 13& 8,8& \\ \hline
$n_\mathrm{id}$ & 1,7& 10,3& 9,9& \\ \hline
$m_\mathrm{F}$ & &1,3& 1,6& 1,8 \\ \hline
$I_0$ & &$4,5\cdot10^{-8}$& $1,3\cdot10^{-7}$& $1,5\cdot10^{-6}$ \\ \hline
$I_{0,\mathrm{TAT}}$, відн.од. & &1& 0,14& 0,04 \\ \hline
$U_d$, B & &0,73& 0,44& 0,12 \\ \hline
$R_\mathrm{TAT}$, відн.од. & &1& 0,54& 0,33 \\ \hline
$K_\mathrm{RECT}$ ($V=0,5$~В)&770 &0,22& 1,33& 5,4 \\ \hline
\end{tabular}
\end{table}


Зростання послідовного опору (приблизно на порядок) призвело до зменшення падіння напруги на діелектричному прошарку.
Як наслідок, напруження електричного поля перестало бути достатнім для ефективного тунелювання Фаулера--Нордгейма і відповідна компонента струму не спостерігається.

При малих зміщеннях ($V<1$~В), як видно з рис.~\ref{figIV_MIS},в, пряма ВАХ опромінених структур добре описується показниковою залежністю
\begin{equation}\label{eqVIsclc}
  I=I_0\,V^{\,m_\mathrm{F}},
\end{equation}
де $m_\mathrm{F}=\frac{V}{I}\frac{\partial V}{\partial I}$ --- диференційний показник ступеня.
Залежність (\ref{eqVIsclc}) характерна для проходження струму, обмеженого просторовим зарядом \cite{SCLC:MA2016,Jafar,SCLC:Kaya}, при цьому величина $m_\mathrm{F}$ відображає енергетичний розподіл пасток, які емітують носії струму.
Для структур після $\gamma$--опромінення і перед УЗО $m_\mathrm{F}\approx1,3$, що свідчить про експоненційний розподіл рівнів пасток по енергії.
Як відомо \cite{SCLC:MA2016,Jafar,SCLC:Kaya}, при цьому $I_0$ залежить від загальної концентрації пасток $N_t$
\begin{equation}\label{eqIoSCLC}
  I_0\sim 1/N_t^{m_\mathrm{F}-1},
\end{equation}
а температурна залежність показника ступеня описується виразом
\begin{equation}\label{eqMT_IoSCLC}
  m_\mathrm{F}=1+T_c/T,
\end{equation}
де
$T_c$ --- характеризує розподіл пасток по енергії:
концентрація рівнів на одиницю енергії $E$ пропорційна $\exp(-E/kT_c)$.

Вираз для струму в SCLC--режимі також часто записують у вигляді \cite{Jafar}:
\begin{equation}\label{eqVIsclcT}
  I(V,T)=C\exp\left(-\frac{E_x}{kT}\right)\,V^{\,m_\mathrm{F}(T)},
\end{equation}
де
$C$ --- певна константа,
$E_x$ --- енергія активації, що визначається положенням заповнених пасток стосовно зони провідності.
На рис.~\ref{figEa_MIS} представлена температурна залежність прямого струму при $V=0,4$~В, скорегована з врахуванням
температурної залежності показника ступеня (\ref{eqMT_IoSCLC}).
Видно, що вона дійсно добре описується виразом (\ref{eqVIsclcT}), причому $E_x=(0,32\pm0,01)$~еВ.

\begin{figure}
\center
\includegraphics[width=0.55\textwidth]{figEa_MIS}%
\caption{\label{figEa_MIS}
Температурна залежність SCLC--струму ($V=0,4$~В)
$\gamma$--опроміненої МОН структури
%Au--SiO$_2$--Si
до УЗО.
Точки --- експеримент,
пряма --- лінійна апроксимація
}%
\end{figure}

Як відомо з літератури \cite{PersenkovBook}, при $\gamma$--опроміненні в системі Si--SiO$_2$
відбувається релаксація механічних напруг,
утворюються нові заряджені дефекти та заряджаються пастки, що вже були раніше.
Внаслідок цих процесів рельєф густини об'ємного заряду вирівнюється.
Вважається, що основними пастками, які захоплюють негативний заряд є центри на інтерфейсній межі розділу,
тоді як позитивний заряд накопичується на $E'$--центрах, $E'_\delta$-- центрах та трьох рівнях, пов'язаних
з напруженими зв'язками \cite{SiO2:Devine,SiO2:Lenahan}.
Проте при фотонному дефектоутворенні (використанні рентгенівських променів чи $\gamma$--квантів) останні рівні утворюються у значно меншій кількості, ніж два перших \cite{SiO2:Devine}.
Зауважимо, для досліджених зразків достатньо низькі температура та парціальний тиск кисню призводили до утворення тонких шарів окису кремнію, проте в літературі \cite{SiO2:Cantin} показано, що при цьому утворюються такі самі радіаційні дефекти, як і для прошарків більшої товщини.

Зазначимо, що ключовим фактором створення електричноактивних радіаційних дефектів у кремнієвих МОН структурах вважається вміст водню \cite{SiO2:Cantin}.
Плівки SiO$_2$ отримані термічним шляхом містять значну кількість атомів водню \cite{PersenkovBook}.
Утворення інтерфейсних рівнів пов'язується з розривом зв'язків $\equiv\!\text{Si}\!-\!\text{H}$ на межі Si--SiO$_2$ за наявності гарячих носіїв заряду \cite{SiO2:Mahapatra,SiO2:Esseni}.
Утворені ненасичені зв'язки $\equiv\text{Si}-$  діють як електронні пастки.
Загалом розрізняють дефекти подібного типу залежно від орієнтації кремнієвої підкладки, на якій вирощено шар окису.
Вважається, що на межі з кремнієм, орієнтованим у площині (111), з'являються $P_b$--центри, для (100) підкладки характерним є поява $P_{b1}$-- та $P_{b0}$--центрів \cite{SiO2:Rev}.
Ці центри хімічно ідентичні, проте характеризуються різною електричною активністю.
Крім того, при $\gamma$--опроміненні систем Si--SiO$_2$, вирощених на (111) підкладках, густина радіаційних поверхневих центрів вища ніж для (100) і досягає величини близько $10^{12}$~еВ/см для дози $10^{7}$~рад \cite{PersenkovBook}.

Відомо що при перевищенні дозою опромінення $\gamma$--квантами величини близько $5\cdot10^5$~рад у структурах Si--SiO$_2$ на основі кремнію з електронною провідністю спостерігається немонотонних розподіл інтерфейсних станів по енергії \cite{PersenkovBook}.
Згідно з даними роботи \cite{ParchSiO2} найбільша густина поверхневих станів спостерігається при
$E_c-(0,32\pm0,04)$~еВ, що збігається з отриманим значенням $E_x$,
Це свідчить, що проходження струму при малих прямих зміщеннях в досліджених структурах після опромінення зумовлене саме захопленням електронів $P_b$--центрами.
Накопичення від'ємного заряду на інтерфейсі також викликає збільшення висоти бар'єру,
що відображається у виявленому зменшенні ТЕ струму.
У літературі і раніше повідомлялося про зменшення струму внаслідок опромінення МОН--структур \cite{SiO2:Niu}.
Проте, як вже зазначалося, основною причиною зменшення ТЕ компоненти струму є дефектоутворення у приповерхневому шарі кремнію та відповідне зростання $n_\mathrm{id}$.

У результаті УЗО, як видно з рис.~\ref{figIV_MIS}, спостерігається збільшення струму, обмеженого просторовим зарядом.
Зокрема, при $t_\mathtt{UST}=60$~хв, цей струм перевищив термоемісійний у всьому дослідженому діапазоні напруг.
Параметри, отримані шляхом апроксимацій ВАХ, зведені в табл.~\ref{tabMIS}.
Виявлене збільшення $I_0$ свідчить, відповідно до формули~(\ref{eqIoSCLC}), про АІ зменшення концентрації $P_b$--центрів, а зростання $m_\mathrm{F}$ (або, що те саме, $T_c$) --- про звуження енергетичного розподілу відповідних рівнів.
З літератури \cite{SiO2:Takakura,SiO2:Wurzer} відомо, що відпал $P_b$--центрів відбувається внаслідок пасивації ненасичених зв'язків атомами водню.
Тобто отримані результати свідчать, що УЗО стимулює дифузію водню.
Загалом подібні процеси спостерігалися і раніше \cite{Ostap:SiO2,Ostap:PhotoLum,ostapenko1997}, проте в нашому випадку мова йде про АІ відпал радіаційних дефектів.


Як відомо \cite{Jafar}, випадок $m_\mathrm{F}=2$ відповідає наявності пасток з однаковим енергетичним рівнем.
Виявлене звуження енергетичного розподілу свідчить про певну вибірковість АІ процесів пасивації:
внаслідок УЗО атоми Н осідають переважно на зв'язки, яким відповідає цілком певне розташування рівнів у забороненій зоні.
На погляд автора, визначальним фактором того, які саме зв'язки будуть заповнені під час АІ відпалу є механічні напруги, неоднорідним чином розподілені на інтерфейсі.
Відомо, що дифузійна здатність домішок у напівпровідниках залежить від механічних напруг \cite{AZIZ2001}
(що може бути основною причиною переміщення водню в полі УЗ).
Відтак, від рівня деформації залежить ефективність АІ пасивації, а УЗО сприяє підвищенню однорідності структури.

Наслідком зменшення від'ємного заряду, накопиченого $P_b$--центрами, є часткове відновлення висоти бар'єру, що викликає зростання ТЕ струму --- див. дані для $I_s$ в табл.~\ref{tabMIS}.
Виявлене зменшення послідовного опору свідчить про частковий АІ відпал радіаційних дефектів у об'ємі напівпровідника.

Звільнені під час утворення $P_b$--центрів при $\gamma$--опроміненні атоми водню можуть
а)~взаємодіяти з іншими зв'язаними атомами Н на межі, викликаючи їхній відрив та утворення додаткових $P_b$--центрів \cite{SiO2:Devine};
б)~переміщуючись у кремнієву підкладку генерувати дефекти та деактивувати легуючу домішку \cite{SiO2:DiMaria};
б)~рухатися в об'єм діелектрика, де, за наявності гарячих носіїв заряду, призводити до появи $E'$--центрів (так званих, об'ємних пасток, bulk--oxide trap) внаслідок розриву зв'язків у системі $\equiv\!\text{Si}\!-\!\text{O}$ \cite{SiO2:Mahapatra,SiO2:Esseni}.
За своєю структурою $E'$--центри --- це вакансії кисню в SiO$_2$ \cite{SiO2:Takakura,SiO2:Devine},
які накопичують позитивний заряд.
Загальна концентрація цих дефектів при $\gamma$--опроміненні з дозою $10^{7}$~рад становить приблизно $10^{18}$~см$^{-3}$, проте по товщині оксидної плівки вони розподілені нерівномірно, найбільша кількість спостерігається біля межі з кремнієм \cite{PersenkovBook}.
Ці центри стабільні при кімнатній температурі \cite{SiO2:Mahapatra}, температура відпалу становить (200$\div$400)$^\circ$C.

Наявність розірваних зв'язків у об'ємі оксиду є причиною появи так званих напруго--індукованих струмів втрат (stress-induced leakage current, SILC) \cite{SiO2:Mahapatra,SiO2:DiMaria}.
 Саме цей струм реєструється в опромінених структурах при зворотному зміщенні (рис.~\ref{figIV_MIS},а та б).
Дійсно, відомо \cite{SiO2:Esseni,SiO2:DiMaria}, що SILC зумовлений тунелювання по пасткам (trap--assisted tunneling, ТАТ) .
Відповідно до \cite{TAT:Gilmore,TAT:GopalSST,TAT:Gopal}, польова залежність ТАТ--струму має вигляд:
\begin{equation}\label{eqIVTAT}
  I_R=I_{0,\mathrm{TAT}}\,(U_d+V_R)\exp\left(-\frac{R_\mathrm{TAT}}{F_m}\right),
\end{equation}
де
параметри $I_{0,\mathrm{TAT}}$ та $R_\mathrm{TAT}$ не залежать від напруги,
$I_{0,\mathrm{TAT}}$ пропорційний концентрації пасток,
$U_d$ --- висота бар'єру.
На рис.~\ref{figIV_MIS} наведено результати апроксимації експериментально отриманих зворотних гілок ВАХ згідно з формулою~(\ref{eqIVTAT}), що підтверджують припущення щодо природу цього струму.
Відхилення від очікуваної залежності при великих зміщеннях може бути зумовлене впливом послідовного опору.
Дані табл.~\ref{tabMIS} свідчать про низькотемпературний ($\sim$80$^\circ$C) АІ відпал радіаційно утворених пасток та зниження висоти бар'єру.
На нашу думку, відпал $E'$--центрів також зумовлений акустодифузією міжвузлових атомів, у даному випадку  кисню.

Збільшення прямого струму та зменшення зворотного в результаті УЗО викликає збільшення коефіцієнта випрямлення
$K_\mathrm{RECT}$, який суттєво зменшився в результаті $\gamma$--опромінення.
У табл.~\ref{tabMIS} наведено значення цієї величини при напрузі 0,5 В.
Дані свідчить, що УЗО при температурах, близьких до кімнатних, здатна частково відновлювати характеристики радіаційно--деградованих  МОН  пристроїв.







\section*{Висновки до розділу \ref{Ch_UST_MW}}
\addcontentsline{toc}{section}{Висновки до розділу \ref{Ch_UST_MW}}
  \begin{enumerate}[leftmargin=0cm,itemindent=3em]
\item Експериментально досліджено вплив мікрохвильового опромінення на параметри точкових дефектів (поперечний переріз захоплення електронів,
розташування енергетичних рівнів у забороненій зоні) в монокристалах $n$--6$H$--SiC, $n$--GaAs та епітаксійних структурах на основі арсеніду ґалію.


\item Показано, що у приповерхневому шарі присутні центри захоплення носії заряду, пов'язані з власними дефектами вакансійного типу.
Виявлені зміни параметрів пасток обумовлені індукованим мікрохвильовим випроміненням
збільшенням кількості міжвузлових атомів.
Показано, що наявність стискуючих напруг у приповерхневому шарі сприяє перетворенню дефектних комплексів внаслідок високочастотного опромінення.

     \item Вперше експериментально досліджено вплив ультразвукової обробки на параметри структур Au--TiB$_x$--$n$--$n^+$--GaAs з контактом Шотткі залежно від частоти та інтенсивності акустичних хвиль.
%Виявлено, що вплив ультразвуку на зворотні гілки ВАХ значно суттєвіший ніж на прямі.
 Експериментально встановлено, що при допороговій інтенсивності (менше 2,5~Вт/см$^2$) характер ультразвукового впливу на зворотний струму залежить від механізму перенесення заряду:
  якщо домінуючим механізмом є тунельний, то УЗО викликає збільшення зворотного струму, якщо термоемісійний --- зменшення.
  Показано, що причиною  виявлених ефектів при допороговій інтенсивності УЗО
  може бути акусто--стимульована дифузія точкових дефектів, тоді як зростання зворотного струму при надпороговій  інтенсивності зумовлене акусто--індукованою генерацією дефектів.
 Виявлено, що зі збільшенням частоти ультразвукової обробки інтенсифікуються процеси перебудови дефектів, які забезпечують проходження тунельного струму.

\item Показано, що ультразвукова обробка здатна викликати зменшення розкиду параметрів і підвищення однорідності характеристик напівпровідникових пристроїв, створених в єдиному технологічному процесі.

\item %Досліджено вплив ультразвукової обробки на парметри структур метал--окис--напівпровідник.
   Виявлено можливість низькотемпературного відпалу радіаційних дефектів у системі Si--SiO$_2$.
   Показано, що акусто--індукований відпал зумовлений дифузією міжвузлових атомів.
 Виявлено, що ультразвукова обробка сприяє звуженню енергетичного спектра радіаційно--індукованих пасток
на інтерфейсі системи    Si--SiO$_2$.


\end{enumerate}	



Основні результати даного розділу представлені в роботах \cite{Olikh:SPQEO2003,Olikh:PJE2004,Olikh:PhChOM2005,Olikh:PZTF2006,Gorb2010,
3Tomsk,50IUFFC,9APTTE,ICU2007GA,6DrogGorb}.



%\cite{PersenkovBook}
%При гамма--опроміненні релаксація механічних напруг в системі Si--SiO$_2$:
%частина дірок, що утворилася при іонізації захоплюється на рівні напружених зв'язків, що призводить до їх розриву
%термічно отримані плівки SiO$_2$ поліморфні,
%%містять велику кількість водню
%
%%Е' центри (пастки дірок) розподілені нерівномірно, найбільша концентрація біля границі з Si.
%%загальна концентрація при 1е7 рад ---  близько 1е18 см-3.
%
%тривалентний кремній --- донороподібний центр,
%відповідає за додатній заряд, який накопичується в об'ємі при опроміненні
%
%%при опроміненні Si--SiO$_2$ утворюються нові заряджені дефекти та заряджаються пастки,
%%що вже були,
%%рельєф густини об'ємного заряду вирівнюється
%
%швидкі поверхневі стани розташовані в приповерхневому шарі напівпровідника,
%повільні --- в приграничній області оксиду.
%
%%густина утворених поверхневих станів при гамма--опроміненні найбільша при (111), найменша (100),
%%при 1е7 рад порядку 1е12 еВ/см2
%
%кореляція між густиною швидких поверхневих станів Nss та пасткових центрів в об'ємі SiO$_2$ No
%Nss=3.75t-6 см/еВ No,
%а отже однакова причина утворення дефектів:наявність напружених валентних зв'язків між атомами та їх розрив в результаті іонізації.
%
%%при опроміненні гамма--квантами найбільша густина станів в околі Ес-(0,3-0,4)еВ (для електронного, для діркового не так),
%%причому така нерівномірність спостерігається лише при перевищенні дозою порогу близько 5е5 рад.
%
%швидкість поверхневої рекомбінації визначається густиною поверхневих станів
%s~Nss
%
%в об'ємі напівпровідника час життя змінюється мало при гамма-опроміненні, в приповерхневому шарі зменшується
%
%%\cite{ParchSiO2}
%%часткова компенсація напівпровідникової підкладинки в процесі радіаційного дефектоутворення
%%гамма--кванти, до 2е7 рад, максимум спектра поверхневих станів Ес-(0,32+-0,04) ев, що характерно для високих доз -- див попередню
%%
%%електронне опромінення (n-тип, (100)) --- максимум густини поверхневих станів Ес-0,3 еВ \cite{LaiSiO2}
%
%
%
%%Interface trap (Nit) утворюються внаслідок розриву зв'язків $\equiv\!\text{Si}\!-\!\text{H}$ на інтерфейсі Si--SiO$_2$,
%%в результаті утворюються $\equiv\text{Si}-N_\text{iT}$, які реєструються як Pb центри.
%%Звільнені атоми водню рухаються в об'єм діелектрика.
%%За наявності гарячих носіїв заряду при цьому генеруються bulk--oxide trap Not внаслідок розриву $\equiv\!\text{Si}\!-\!\text{O}$.
%%Розірвані зв'язки $\equiv\!\text{Si}\!-\!\text{O}$ в об'ємі оксиду викликають появу stress-induced leakage current (SILC).
%%However, unlike $\equiv\!\text{Si}\!-\!\text{H}$  bonds, broken $\equiv\!\text{Si}\!-\!\text{O}$ bonds are not known to recover
%%at room temperature after the stress is removed.
%%В роботі пропонується, що $\equiv\!\text{Si}\!-\!\text{H}$ bonds are broken by holes (and not by electrons). \cite{SiO2:Mahapatra}
%
%
%%Гарячі електрони ведуть до звільнення водню з інтерфейсу Si--SiO$_2$,
%%міграція Н в оксиді викликає появу bulk--oxide trap.
%%Механізм SILC --- trap--assisted tunneling
%%The results reported so far demonstrate that in the channel hot electron stress conditions there is no correlation between and SILC generation.
%%В самій роботі Nit generation is evidently related to stress-induced leakage current (H/D) release \cite{SiO2:Esseni}
%%
%%SILCs  (This  phenomenon  is  often  referred  to  as  trap-assisted  tunneling ) can  be  best  explained  by  the generation  of  neutral  electron  traps  in  the  oxide  layer.
%%These  studies  have  directly  shown  the  relationship  of SILCs  to  the  trap  creation  phenomenon.
%%Trapped  holes  or  sites  produced by  trapped-hole  annihilation  by  free  electrons  will  be demonstrated  not  to  contribute  to  SILC  for  most  stressing conditions.
%%Since  the  oxide  deterioration  producing  SILC  will  be shown  here  to  be  caused  by  electron  heating  in  the  oxide layer,
%%oxide  defect  formation  due  to  hot  carriers  will  be briefly  reviewed.
%%The  production  of  interface  states or  traps by  the  annihilation  of  trapped  holes  and/or  their  presence  is
%%one  component  of  the  oxide  degradation  process.
%%The other  component  is  due  to  the  trap  creation  process  also  caused by  hot  electrons.
%%  Also, the  most  significant  component  contributing  to  the  SILC  in the  experiments  discussed here  appears to  be due  to  the  generation  of  the  neutral-electron  traps,  not  the  interface  states and/or  anode positive  charge.
%%Also,  generation--recombination  sites  and boron  deactivation  can be produced  in  the Si  substrate by  the hydrogen.
%%However,  during  trap  creation,  many  different  defect  sites  with  different  processing  dependencies  from  one wafer  run  to  another  can  occur.  These  defects  include  interface  states, generation-recombination  centers  in  the  Si  substrate,  neutral  electron  traps  in  the  oxide  bulk  distributed  away  from  the  cathode/oxide  interface,  and positively  charged  oxide  sites  (mostly  <<slow  states>> )  near the  anode/oxide  interface.\cite{SiO2:DiMaria}
%
%%На границі з (111) --- Pb центр,
%%на границі з (100) --- Pb1 та Pb0.
%%центри хімічно ідентичні, різна електрична активність (Pb1, здається не активний)
%%взагалі робота --- огляд.. там багато всякого про інтерфейсні рівні \cite{SiO2:Rev}
%
%%Опромінення гамма--квантами, доза 1е6,
%%утворилися нові інтерфейсні рівні (порядку 3е10 см-2).. головна ідея -- при прикладенні напруги в опромінених структурах генеруються більше пасток \cite{AASSIME1997}
%
%%тонкі (20--40 ангстрем) шари , опромінення х-променями (10 кев), утворюються такі самі дефекти (Е' (позитивно заряджені вакансії кисню) та два типи Pb, бо (100) пластинки)
%%that one key parameter for the moddication of the electrically  active  defects  in  MOS  structures  is their hydrogen content \cite{SiO2:Cantin}
%
%%опромінення протонами (еквівалентна доза гамма--квантів близько 1е6) --- зміни струмів різних (там транзистор) можуть бути і від'ємні після опроміненя
%%is  shown  to be  a result  of  the increase  in  fixed oxide  charge  and  surface  trap  densities,  in combination with the different SRH  recombination rate limiting
%%mechanisms in the presence of high density surface traps \cite{SiO2:Niu}
%
%%огляд радіаційних дефектів в SiO2:
%%пастки, які захоплюють позитивний заряд--- відпалюються при температурах 200--400 градусів цельсію
%%Е' центр (вакансія кисню) --- основний,
%%$E'_\delta$ (комплекс вакансія кисню--кремній), тяжко утворюється при дії фотонів (гамма та Х)
%%три пастки пов'язані з напруженими зв'язками
%%негативний заряд
%%Pb центри (причому звільнені атоми водню можуть взаємодіяти і іншими, ще зв'язаними, і відривати їх також)
%%як не дивно, знову вакансії кисню
%%we generally expect the neutral oxygen-vacancy to be a hole trap, it may also trap electrons leading  to a negatively charged state in the SiO2  bandgap
%%about 1  eV below the conduction band edge \cite{SiO2:Devine}
%
%%decrease of base current during annealing. We  found that two mechanisms
%%are responsible for annealing: the recombination  of charges and
%%the bonding of hydrogen atoms on interface traps.
%%процеси з двома енергіями активації -- 0,23 ев та 0,046ев
%%As  the  Shockley-Read-Hall  current  is  proportional to the
%%density of  interface traps,  the decrease of  the base current
%%is proportional to  the  density of  the  annealed interface traps. \cite{SiO2:Wurzer}
%
%%Ще одна робота про результат опромінення: переважно Е', Pb (або два варіанти для 100) все повторюється..
%%і те, що Pb амфотерний (може захоплювати і дірки, якщо його рівень в нижній половині забороненої зони)\cite{SiO2:Lenahan}
%
%%It is well established nowadays that one of the main hole trapping centres responsible for irradiation-induced positive
%%charge is the E'-centre. Electron spin resonance (ESR)
%%studies have revealed that the E'-centre is a defect resulting
%%from an oxygen vacancy in SiO2. A direct correlation has been
%%established between the annealing of the oxide-trapped charge.
%%The key SiO2/Si interface defect is called Pb-centre and
%%corresponds to a silicon dangling bond at the interface, back-
%%bonded by three other silicon atoms . Passivation of these dangling bonds (DBs) occurs
%%through binding with hydrogen.
%%From this result it is concluded that the degradation
%%and recovery of the device performance are mainly caused by
%%trapped-oxide charge
%%Вираз для струму, пов'язаного з поверхневою рекомбінацією, близький до того, що використовували ми
%%\cite{SiO2:Takakura}
%
%%залежність дифузійної здатності від механічних напруг \cite{AZIZ2001}
%%
%%\cite{Ostap:SiO2} стан границі Si--SiO2, підсилення дифузії водню
%%\cite{Ostap:PhotoLum,ostapenko1997} підсилення дифузії водню
